% Proposed revision text only. Not included in the manuscript automatically.
\paragraph{Revised digital twin and profiling validation.}
We model process-attributed power and a dimensionless composite QoS score as
functions of the recorded offered load. The revised twin comprises three neural
networks and one nonnegative piecewise-linear regression. DPDK power uses a
16-unit ReLU hidden layer and a softplus output in watts. The two QoS networks
retain the original 64/32/16-unit hidden architecture, with a sigmoid output in
the original score units. USR power uses linear interpolation between seven
nonnegative fitted knot powers. This replaces the historical unconstrained USR
power network, which produced negative predictions at measured-range loads.
All original artifacts are retained for historical comparison.

Complete profiling sweeps 1--3 are used for parameter learning and input-scaler
fitting, sweep 4 for checkpoint and candidate selection, and sweep 5 for
internal evaluation. No sweep is divided between these roles. The five sweeps
come from the same campaign for each implementation and have been inspected
previously; the internal evaluation is therefore neither a new physical campaign
nor a fully blind confirmatory test. Model choices are frozen before workload
diagnostics and do not use controller outcomes. The USR power choice additionally
reflects the substantially smaller low-load errors of the preserved regression
relative to the constrained neural alternative; this is a documented scientific
tradeoff, not a prespecified equivalence test.

\paragraph{Recovered composite score.}
The experiments used four components, including received throughput. The earlier
three-component description was a reporting discrepancy; the historical labels
are preserved. Define $\operatorname{clip}_{[a,b]}(x)=\min(b,\max(a,x))$.
For recorded offered load $L$ in Mbit/s, DN-side received throughput $R$ in
Mbit/s, stored downlink delay $d$ in milliseconds, RAN jitter $j$ in microseconds,
and stored interval packet-loss fraction $\ell$, the score is
\begin{align}
 T &= \operatorname{clip}_{[0,1]}\!\left(\frac{R}{L+10^{-9}}\right), &
 D &= 1-\frac{\operatorname{clip}_{[0,40]}(d)}{40},\\
 J &= 1-\frac{\operatorname{clip}_{[0,10]}(j)}{10}, &
 B &= 1-\frac{\operatorname{clip}_{[0,5]}(100\ell)}{5},\\
 Q &= 0.40T+0.20D+0.20J+0.20B.
\end{align}
The mappings are respectively \texttt{DN\_Rx\_Mbits/s},
\texttt{Objective Throughput (DL/UL)}, \texttt{downlink\_delay\_ms},
\texttt{RAN\_jitter}, and \texttt{RAN\_PLR\_interval}.
The value $Q\geq0.90$ is a composite optimization target. It does not establish
formal 5QI compliance, full throughput delivery, or a per-packet delay guarantee.
The original received-traffic direction and loss-denominator construction retain
the qualifications stated below; no new physical interpretation is inferred
merely from numerical reproduction of the labels.

\paragraph{Actions, outcomes and units.}
The available actions are DPDK, USR(1), and USR(2). For USR with $k$ instances,
we query the USR surrogates at $L/k$, set total attributed power to $kP(L/k)$,
and retain the per-instance score $Q(L/k)$. This is a balanced-load approximation,
without measured multi-instance interference or imbalance. The fitted
interpolation domain is 0.01--5,000 Mbit/s per instance; unmeasured zero load
and out-of-domain inputs are rejected. All three actions share a total-load
domain of 0.02--5,000 Mbit/s, which contains every audited training/validation
workload input. Nonnegative power and bounded QoS follow from the output
parameterizations, without post-hoc prediction clipping. These constraints do
not establish accuracy between measured load levels.

For quarter-hour intervals, modeled energy is
$E=\sum_t P_t\times0.25$ Wh. Slot-specific power divided by offered load has
units W/(Mbit/s), equivalently J/Mbit. Its arithmetic mean differs from aggregate
energy per offered Mbit. Reconfiguration penalties are reward preferences and
are not added to energy as measured startup electricity. Decisions at interval
start may use only previously completed telemetry and causal forecasts; the
realized interval load is used to score the resulting action afterwards.

\paragraph{Limitations and scope of controller comparisons.}
Power targets are Scaphandre process-attributed estimates, not whole-server
electrical measurements. DPDK targets numerically match the sum of the catalogued
process exports, while USR targets match individual process exports. The exact
historical joins, metering uncertainty and idle behavior are not authenticated.
The offered-load objective is labelled per direction; its mapping to aggregate
workload traffic and the physical DN/RAN flow direction are not fully established.
The loss numerator is supported by counter differences, but the interval-loss
denominator and reset/aggregation handling remain unresolved.

QoS networks estimate a conditional point value from offered load. Severe
measured outcomes below 0.90 remain among their predicted-safe observations.
Most DPDK score variation occurs between measurements at identical load, so
larger load-only networks cannot be assumed to recover the omitted state.
Empirical sweep tails and residual sensitivity scenarios are descriptive; they
are not calibrated operational probabilities, confidence intervals or safety
guards. The fixed packet size, exposed same-campaign samples and balanced
USR(2) approximation limit physical generalization. Controller experiments using
this twin compare algorithms conditional on these models and assumptions. They
cannot establish deployed QoS guarantees, service-continuous switching or
whole-server energy savings.

% Suggested table caption:
% Grouped profiling prediction errors of the four frozen revision surrogates.
% V denotes sweep 4; T denotes the exposed same-campaign internal sweep 5.
% Power errors are watts per instance; QoS errors are dimensionless composite
% score units. Bias is prediction minus measurement. These are descriptive
% within-campaign estimates, not independent physical certification.
%
% Suggested figure captions:
% (1) Predictions versus profiling measurements, separately for validation and
% internal evaluation. The diagonal indicates equality; overlap hides repeated
% observations at identical load.
% (2) Mean residual and empirical 5th--95th residual percentiles by offered load.
% Shading describes correlated measurements and is not a confidence interval.
% (3) Empirical below-target fractions and lower score tails across five complete
% sweeps. These are observed same-campaign frequencies, not calibrated future risk.
% (4) Synthetic-load power rankings under frozen and retained alternative USR
% power models. Differences are pointwise surrogate sensitivity, not controller
% energy savings or physical configuration-ranking validation.
